\chapter{Discussion and Conclusion}
\label{chap:discussion}

At the start of this thesis, we asked whether an LLM can generate pedagogically sound Productive Failure tasks that come close to the tasks designed by PF experts, rather than polished story problems that students can already solve with one taught procedure. Because the surface rubric from prior work could not tell these two kinds of task apart, answering this question also required a way to check the PF requirements of a generated task automatically. In this chapter, we first revisit our four research goals (Section~\ref{sec:disc-goals}) and then discuss what the results imply for the use of LLMs in education more broadly (Section~\ref{sec:disc-findings}), and the limitations of our evidence together with directions for future work (Section~\ref{sec:disc-limitations}), before we conclude (Section~\ref{sec:disc-conclusion}).

\section{Revisiting the Research Goals}
\label{sec:disc-goals}

The four goals of this thesis build on each other. G1 examined where one-shot generation and its rubric-based evaluation fall short, G2 and G3 addressed these two shortcomings with a staged generation pipeline and a requirement-based evaluation, and G4 tested whether the pipeline still works across many topics, where the generator must also identify what students already know. In the following, we summarize what Chapters~\ref{chap:oneshot} to~\ref{chap:results} found for each goal and then relate the results back to the expert-designed tasks that motivated this work.

\textbf{G1: One-Shot Generation and Rubric Evaluation}\\[0.2em]
Our first goal was to reproduce one-shot LLM generation of PF tasks from prior work and to identify where both the generated tasks and their rubric-based evaluation fall short. In Chapter~\ref{chap:oneshot} and Section~\ref{sec:res-oneshot}, we found that one-shot prompting writes tasks that look like Productive Failure but rarely work like it. GPT-5 mini scores $14.6$ out of 15 on the rubric of Rahimi et al.~\cite{rahimi2025}, above the $13.38$ they report for GPT-4o, yet familiar methods fall short on only 12 of its 118 tasks ($10\%$), and GPT-4.1 does not do better with 13 of 120 ($11\%$). The rubric misses this failure because it rates the story and wording of a task rather than testing what its numbers allow students to do.

\textbf{G2: Staged Generation with \textsc{StructPF}}\\[0.2em]
To address the weaknesses of the generated tasks, our second goal was to design \textsc{StructPF}, a pipeline that generates PF tasks in smaller, checked steps instead of one pass. Study~A in Chapter~\ref{chap:results} shows that this staged design makes familiar methods fall short far more often. Because Stage~2 plans which familiar shortcut should fail and verifies the task numbers with calculations before Stage~3 writes the text, familiar methods fall short on 96 of 120 \textsc{StructPF} tasks ($80\%$), compared with $10\%$ for one-shot GPT-5 mini given the same prior knowledge (Section~\ref{sec:res-progression}). This gain does not come at the cost of task soundness, since every \textsc{StructPF} task stays on target, $95\%$ are mathematically valid, and room for invention stays at the one-shot level ($73\%$ versus $72\%$).

\textbf{G3: Requirement-Based Evaluation with \textsc{PF Eval}}\\[0.2em]
To address the weaknesses of the rubric, our third goal was to develop \textsc{PF Eval}, an evaluation that tests the three PF requirements directly, and to validate it on expert-designed, control, and human-rated tasks. The validation in Chapter~\ref{chap:eval-framework} shows that testing the requirements directly separates working PF tasks from flawed ones, which the surface rubric could not do. On 29 reference tasks with known strengths and flaws, \textsc{PF Eval} matches every expected outcome, accepting the published expert tasks and rejecting each routine, one-shot, and control task on the measure it breaks (Section~\ref{subsec:val-known-tasks}). Its prior-knowledge solver also recovers $85.5\%$ of 256 expected student approaches and, on every published PF task, tries at least one familiar method that falls short (Section~\ref{subsec:val-solver}). On a blinded 24-task sample from Study~A, \textsc{PF Eval} also agrees with human ratings on 22 of the 24 tasks overall (Section~\ref{sec:res-human-eval}).

\textbf{G4: Comparison across Many Mathematics Topics}\\[0.2em]
Our final goal was to compare \textsc{StructPF} with one-shot generation across many mathematics topics rather than the single basketball setting tested in prior work~\cite{rahimi2025}. In Study~B of Chapter~\ref{chap:results}, the generator receives only a target concept from 30 Common Core topics in Grades~5 to~10 and has to identify what students already know on its own, and \textsc{StructPF} with grounding produces the most working tasks in this setting. \textsc{StructPF} with grounding passes both steps of the Two-Step Check on 19 topics, compared with 16 for \textsc{StructPF} without grounding and 6 for the one-shot prompt (Section~\ref{sec:res-grounding}). Much of this gap arises before the task is even tested, because the one-shot prompt lists the target concept itself as prior knowledge on 17 of the 30 topics and \textsc{StructPF} without grounding on 5, whereas grounding on the Learning Commons graph removes these target leaks entirely.

\textbf{Overall Distance to Expert Tasks}\\[0.2em]
Taken together, these results show that \textsc{StructPF} closes much of the gap between one-shot generation and expert-designed PF tasks, but does not yet reach expert level. Under the same five measures of \textsc{PF Eval}, 11 of the 12 published expert tasks pass every measure on their own topics, compared with 68 of 120 \textsc{StructPF} tasks ($57\%$) and only 10 of 118 one-shot GPT-5 mini tasks ($8.5\%$) in Study~A. Most of the remaining gap comes from tasks that give away the solution route or still let a familiar method settle the question (Section~\ref{sec:res-progression}). Although our blinded human verification confirms that \textsc{PF Eval} closely follows human judgment on these structural measures, the results still show how closely the generated tasks follow expert design rather than whether real students learn from them, which we return to in Section~\ref{sec:disc-limitations}.

\section{Implications for LLMs in Education}
\label{sec:disc-findings}

Although this thesis studies Productive Failure, the difficulties we encountered reflect a broader tension in using LLMs for teaching. LLMs hold vast subject knowledge, but teaching also depends on how learners think before instruction, and this knowledge is scarce both in the training data and in the way these models generate text. We summarize four lessons that apply beyond PF task design and then describe how our approach addresses them.

\textbf{Missing Knowledge about Learners}\\[0.2em]
Larger models improve on what is well documented in expert writing, but not on what learners do. GPT-5 mini clearly outperforms GPT-4.1 on mathematical validity ($92\%$ versus $57\%$), yet familiar methods fall short on almost the same share of their tasks ($10.2\%$ versus $10.8\%$, Section~\ref{sec:res-progression}). Most mathematical text records the end point of learning, while the reasoning of younger students before instruction and the links between knowledge components across grades are rarely written down. This is also why the one-shot prompt listed the target concept as prior knowledge on 17 of 30 topics, and why only curriculum grounding removed these leaks (Section~\ref{sec:res-grounding}). Knowledge about learners therefore has to come from curricula, teachers, and classroom data rather than from a larger model.

\textbf{Bias toward the Most Common Solution}\\[0.2em]
Because a language model favors probable continuations, it concentrates on the most frequent solution, while real student thinking is spread over many approaches, several of them rare. On Kapur and Bielaczyc's bicycle problem, 19 of 24 simulated attempts used single-variable algebra and none reproduced the four other pathways of real Grade~7 groups (Section~\ref{subsec:res-kapur-pathways}). Recent studies likewise find that simulated students outperform real ones and abandon their assigned errors~\cite{srivatsa2025abilities,do2026faithfulness,scarlatos2026simstudents}. The same tendency affects generation, since numbers on which familiar methods fall short are an unlikely choice when a model writes a story freely, and only planning and verifying them in Stage~2 raised this share from $10\%$ to $80\%$. Simulated students can therefore test what is solvable with given skills, as the solver in \textsc{PF Eval} does, but they cannot stand in for a classroom.

\textbf{Judges That Reward Polish}\\[0.2em]
An LLM judge that reads only the task text shares the blind spots of the generator and tends to reward polish~\cite{zheng2023judging}. Our rubric therefore rated one-shot GPT-5 mini above \textsc{StructPF} ($14.6$ versus $12.1$ of 15), although familiar methods fell short on far fewer of its tasks. Whenever LLMs generate and judge teaching material, evaluation should test what the material makes a learner do, and the judge should first be validated on material with known strengths and flaws, as we did for \textsc{PF Eval}.

\textbf{Giving Away the Answer}\\[0.2em]
LLMs are trained to be helpful and to resolve the question, whereas PF deliberately leaves the key step to the learner. Even \textsc{StructPF} stated the solution route in 32 of 120 tasks although its prompts forbid this (Section~\ref{sec:res-progression}). Restraint therefore cannot simply be requested in the prompt, but has to be built into the process and checked on the output, as Puech et al.~\cite{puech2025} also had to do for an LLM tutor in PF.

\textbf{Our Approach for Generation and Evaluation}\\[0.2em]
These problems do not make LLMs unsuitable for pedagogical work, but they change how we should use them. Our experience with \textsc{StructPF} and \textsc{PF Eval} suggests four principles that apply equally to other teaching material, such as hints, quiz questions, or tutoring dialogues.

First, the process has to follow the pedagogical intention behind the material. \textsc{StructPF} does not split generation into steps for its own sake. Each stage corresponds to a decision that a PF designer makes for a specific purpose, such as choosing numbers so that a familiar method falls short, and giving this decision its own step accounts for most of the gain over one-shot generation.

Second, knowledge about how students think has to come from experts and teachers. An LLM cannot think like a student, and a larger model does not close this gap. In our pipeline, this knowledge enters through the curriculum graph, which removed all target leaks, and through published expert tasks and classroom pathways, which define the expected outcomes in \textsc{PF Eval}.

Third, we have to respect the limits of LLMs in education. Although these models hold broad subject knowledge, they are distribution models that favor the most likely expert continuation rather than reasoning like a novice learner. With this limit in mind, we designed the prior-knowledge solver in \textsc{PF Eval} to do its best using only the listed skills instead of relying on simulated students to role-play mistakes, and this design worked reliably across our validation tasks.

Fourth, every output has to be verified. LLMs write in a confident tone that suggests complete knowledge, yet they still make obvious mistakes. In Stage~2, for example, the model proposed a right triangle with legs 3 and 4 but a hypotenuse of 6 (Table~\ref{tab:d2b-number-regimes}). We therefore check every numerical condition with a calculator and verify quoted task evidence in code, rather than trusting how convincing the text reads.

\section{Limitations and Future Work}
\label{sec:disc-limitations}

Our evidence comes from automated evaluation and a 24-task blinded human verification rather than classroom studies. The following four limitations bound what our results can show, and each points to a direction for future work.

\subsection{No Evaluation with Real Students}
The most important limitation is that no generated task has been tried with real students. While our blinded human review confirms the structural PF measures on a sample of tasks, all large-scale outcomes come from model judges and a solver restricted to prior knowledge, so they show that a task is designed like a PF task, not that students actually invent several approaches, fail productively, or learn more in the instruction phase that follows. Simulated students cannot fill this gap, since they did not reproduce the approaches of real classrooms (Section~\ref{sec:res-negative-signals}). Future work should therefore test the generated tasks in classrooms with multiple independent teacher raters and real students, comparing them with expert tasks on the approaches that students produce and on what they learn after instruction.

\subsection{Few Expert-Written PF Tasks}
Only 12 published expert PF tasks were available to us~\cite{puech2025}. They serve both as the positive references for validating \textsc{PF Eval} and as the expert level in our comparison, but they cover only a few topics. Section~\ref{sec:disc-goals} therefore compares expert tasks on their own topics with generated tasks on other topics, and a fair comparison would require experts to write new PF tasks for the same topics as the generator. Future work should collect such topic-matched expert tasks, which would also enlarge the reference set for validating \textsc{PF Eval}.

\subsection{Subject and Curriculum Scope}
Our studies cover only mathematics in English and follow the US Common Core for Grades~5 to~10 through the Learning Commons graph, with 30 tasks per setting in Study~A and one task per topic in Study~B. Whether the same pipeline works for other subjects, curricula, and languages remains untested. Future work could apply \textsc{StructPF} to other subjects such as science and to other curricula by replacing the Learning Commons graph with the corresponding curriculum data.

\subsection{Choice of Generator Model}
Due to our budget, we generated all \textsc{StructPF} tasks with GPT-5 mini rather than with larger frontier models such as GPT-5 or Claude Opus~5.5. As shown in Table~\ref{tab:disc-token-cost}, because \textsc{StructPF} makes several calls per task, generating 120 tasks costs about $\$7$ with GPT-5 mini, compared with roughly $\$35$ with GPT-5 and $\$75$ with Claude Opus~5.5 before iterative development and evaluation. Larger models might follow the design rules more reliably and reduce the remaining tasks that give away the solution route or that a familiar method still solves. However, the comparison of GPT-5 mini and GPT-4.1 in Section~\ref{sec:disc-findings} suggests that a stronger model improves subject knowledge rather than knowledge of how students think, so we expect this gap to remain regardless of the model. Future work should rerun \textsc{StructPF} with these models to test which failures disappear and which remain.

\begin{table}[!htb]
\centering
\caption[Calls, Tokens, and Estimated API Cost per Task]{Calls, Tokens, and Estimated API Cost per Task}
\label{tab:disc-token-cost}
\vspace{0.3em}
\footnotesize
\setlength{\tabcolsep}{5.5pt}
\renewcommand{\arraystretch}{1.2}
\begin{tabular}{@{} l l c r r r r r @{}}
\toprule
\textbf{Workflow} & \textbf{Model} & \textbf{Calls} & \textbf{Input} & \textbf{Output} & \textbf{Total} & \textbf{Cost / Task} & \textbf{Cost / 120} \\
\midrule
One-Shot Gen & GPT-5 mini & $1.0$ & $1.6\text{k}$ & $2.0\text{k}$ & $3.6\text{k}$ & $\$0.004$ & $\$0.50$ \\
One-Shot Gen & GPT-4.1 & $1.0$ & $1.6\text{k}$ & $0.3\text{k}$ & $1.9\text{k}$ & $\$0.006$ & $\$0.70$ \\
\textsc{StructPF} Gen & GPT-5 mini & $8.0$ & $36.0\text{k}$ & $24.0\text{k}$ & $60.0\text{k}$ & $\$0.060$ & $\$7.00$ \\
\textsc{PF Eval} Round & GPT 5.6 Luna + GPT-4.1 & $4.5$ & $19.0\text{k}$ & $9.5\text{k}$ & $28.5\text{k}$ & $\$0.024$ & $\$2.90$ \\
\midrule
\textsc{StructPF} Gen (est.) & GPT-5 & $8.0$ & $36.0\text{k}$ & $24.0\text{k}$ & $60.0\text{k}$ & $\$0.290$ & $\$35.00$ \\
\textsc{StructPF} Gen (est.) & Claude Opus 5.5 & $8.0$ & $36.0\text{k}$ & $24.0\text{k}$ & $60.0\text{k}$ & $\$0.620$ & $\$75.00$ \\
\bottomrule
\end{tabular}
\end{table}

\section{Conclusion}
\label{sec:disc-conclusion}

In this thesis, we asked whether an LLM can generate Productive Failure tasks that come close to expert design. One-shot generation produces polished tasks that a rubric rates highly, yet familiar methods fall short on only about one in ten of them. We therefore based both generation and evaluation on how a PF expert designs and checks a task. \textsc{StructPF} follows the decisions of a PF designer, takes prior knowledge from curriculum data, and verifies the numbers before the task is written. \textsc{PF Eval} checks the PF requirements on the finished task alone, with a solver restricted to prior knowledge, matching all 29 expected outcomes on our reference tasks and closely following blinded human ratings on generated tasks. Familiar methods fall short on $80\%$ of \textsc{StructPF} tasks compared with $10\%$ of one-shot tasks, and $57\%$ compared with $8.5\%$ pass all five measures, although \textsc{StructPF} does not yet reach the published expert tasks and has not been tested with real students.

Beyond Productive Failure, our results show that LLMs bring broad subject knowledge and fluent writing to education, yet they remain distribution models that favor the most likely expert continuation rather than thinking like a learner. Using them effectively in education therefore means respecting these limits, as our prior-knowledge solver does instead of role-play simulated students, supplying knowledge about learners from experts, teachers, and curricula, and verifying every output rather than trusting its confident tone. We hope that \textsc{StructPF} and \textsc{PF Eval} offer a starting point for using LLMs in this way, so that the tasks they generate let students struggle with the key idea before they are taught it.
